% GENERATED FILE. Edit canonical lessons, then run npm run generate:books.
% canonical-source-hash: fnv1a64:02efa6d281d718a9
% canonical-lessons: TE-C246-appu-ivvu, TE-C246-tirigi-ivvu, TE-C246-khali-ceyu, TE-C246-taggu, TE-C246-nirakarincu

\chapter{To lend, To give back, To empty, To decrease, To refuse}
\label{ch:te-to-lend-to-give-back-to-empty-to-decrease-to-refuse}
\hlchaptermodality{246}

\begin{chapteropening}
\textbf{By the end of this chapter:} \emph{I can say to lend, to give back, to empty, to decrease and to refuse.}

Complete the last lesson of chapter 246: I can say to lend, to give back, to empty, to decrease and to refuse.
\end{chapteropening}

\section[\texorpdfstring{\te{అప్పు} \te{ఇవ్వు} (appu ivvu)}{అప్పు ఇవ్వు (appu ivvu)}]{\te{అప్పు} \te{ఇవ్వు} (appu ivvu) — to lend (money)}

\label{lesson:TE-C246-appu-ivvu}



\begin{quote}
Before the new one: say the Telugu for to download, then the Telugu for to type.
\end{quote}

\subsection*{You'll want to know: \te{అప్పు} \te{ఇవ్వు}}
\textbf{\te{అప్పు} \te{ఇవ్వు}} — \emph{appu ivvu} — \textquotedblleft{}to lend (money)\textquotedblright{}.

\begin{grammarlens}[title={What you've built}]
One more word for this chapter.
\end{grammarlens}

\subsection*{Your turn}
\begin{itemize}
\raggedright
  \item \emph{Say it:} \emph{appu ivvu}
  \item \emph{Say it:} \emph{appu ivvu}, once more
  \item \emph{Say it:} say \emph{appu ivvu}
  \item \emph{Recall:} say the Telugu for to download, then the Telugu for to type, then say \emph{appu ivvu} again
\end{itemize}

\begin{tcolorbox}[breakable,colback=teal!4,colframe=teal!35!black,title={Before you move on}]
What does \te{అప్పు} \te{ఇవ్వు} mean? (\textquotedblleft{}To lend (money)\textquotedblright{}.) Say it once more.
\end{tcolorbox}
\section[\texorpdfstring{\te{తిరిగి} \te{ఇవ్వు} (tirigi ivvu)}{తిరిగి ఇవ్వు (tirigi ivvu)}]{\te{తిరిగి} \te{ఇవ్వు} (tirigi ivvu) — to give back, to return}

\label{lesson:TE-C246-tirigi-ivvu}



\begin{quote}
Before the new one: say the Telugu for to type, then the Telugu for to lend.
\end{quote}

\subsection*{You'll want to know: \te{తిరిగి} \te{ఇవ్వు}}
\textbf{\te{తిరిగి} \te{ఇవ్వు}} — \emph{tirigi ivvu} — \textquotedblleft{}to give back, to return\textquotedblright{}.

\begin{grammarlens}[title={What you've built}]
One more word for this chapter.
\end{grammarlens}

\subsection*{Your turn}
\begin{itemize}
\raggedright
  \item \emph{Say it:} \emph{tirigi ivvu}
  \item \emph{Say it:} \emph{tirigi ivvu}, once more
  \item \emph{Say it:} say \emph{tirigi ivvu}
  \item \emph{Recall:} say the Telugu for to type, then the Telugu for to lend, then say \emph{tirigi ivvu} again
\end{itemize}

\begin{tcolorbox}[breakable,colback=teal!4,colframe=teal!35!black,title={Before you move on}]
What does \te{తిరిగి} \te{ఇవ్వు} mean? (\textquotedblleft{}To give back, to return\textquotedblright{}.) Say it once more.
\end{tcolorbox}
\section[\texorpdfstring{\te{ఖాళీ} \te{చేయు} (khāḷī cēyu)}{ఖాళీ చేయు (khāḷī cēyu)}]{\te{ఖాళీ} \te{చేయు} (khāḷī cēyu) — to empty}

\label{lesson:TE-C246-khali-ceyu}



\begin{quote}
Before the new one: say the Telugu for to lend, then the Telugu for to give back.
\end{quote}

\subsection*{You'll want to know: \te{ఖాళీ} \te{చేయు}}
\textbf{\te{ఖాళీ} \te{చేయు}} — \emph{khāḷī cēyu} — \textquotedblleft{}to empty\textquotedblright{}.

\begin{grammarlens}[title={What you've built}]
One more word for this chapter.
\end{grammarlens}

\subsection*{Your turn}
\begin{itemize}
\raggedright
  \item \emph{Say it:} \emph{khāḷī cēyu}
  \item \emph{Say it:} \emph{khāḷī cēyu}, once more
  \item \emph{Say it:} say \emph{khāḷī cēyu}
  \item \emph{Recall:} say the Telugu for to lend, then the Telugu for to give back, then say \emph{khāḷī cēyu} again
\end{itemize}

\begin{tcolorbox}[breakable,colback=teal!4,colframe=teal!35!black,title={Before you move on}]
What does \te{ఖాళీ} \te{చేయు} mean? (\textquotedblleft{}To empty\textquotedblright{}.) Say it once more.
\end{tcolorbox}
\section[\texorpdfstring{\te{తగ్గు} (taggu)}{తగ్గు (taggu)}]{\te{తగ్గు} (taggu) — to decrease, to go down}

\label{lesson:TE-C246-taggu}



\begin{quote}
Before the new one: say the Telugu for to give back, then the Telugu for to empty.
\end{quote}

\subsection*{You'll want to know: \te{తగ్గు}}
\textbf{\te{తగ్గు}} — \emph{taggu} — \textquotedblleft{}to decrease, to go down\textquotedblright{}.

\begin{grammarlens}[title={What you've built}]
One more word for this chapter.
\end{grammarlens}

\subsection*{Your turn}
\begin{itemize}
\raggedright
  \item \emph{Say it:} \emph{taggu}
  \item \emph{Say it:} \emph{taggu}, once more
  \item \emph{Say it:} say \emph{taggu}
  \item \emph{Recall:} say the Telugu for to give back, then the Telugu for to empty, then say \emph{taggu} again
\end{itemize}

\begin{tcolorbox}[breakable,colback=teal!4,colframe=teal!35!black,title={Before you move on}]
What does \te{తగ్గు} mean? (\textquotedblleft{}To decrease, to go down\textquotedblright{}.) Say it once more.
\end{tcolorbox}
\section[\texorpdfstring{\te{నిరాకరించు} (nirākariñcu)}{నిరాకరించు (nirākariñcu)}]{\te{నిరాకరించు} (nirākariñcu) — to refuse}

\label{lesson:TE-C246-nirakarincu}



\begin{quote}
Before the new one: say the Telugu for to empty, then the Telugu for to decrease.
\end{quote}

\subsection*{You'll want to know: \te{నిరాకరించు}}
\textbf{\te{నిరాకరించు}} — \emph{nirākariñcu} — \textquotedblleft{}to refuse\textquotedblright{}.

\begin{grammarlens}[title={What you've built}]
One more word for this chapter.
\end{grammarlens}

\subsection*{Your turn}
\begin{itemize}
\raggedright
  \item \emph{Say it:} \emph{nirākariñcu}
  \item \emph{Say it:} \emph{nirākariñcu}, once more
  \item \emph{Say it:} say \emph{nirākariñcu}
  \item \emph{Recall:} say the Telugu for to empty, then the Telugu for to decrease, then say \emph{nirākariñcu} again
\end{itemize}

\begin{tcolorbox}[breakable,colback=teal!4,colframe=teal!35!black,title={Before you move on}]
What does \te{నిరాకరించు} mean? (\textquotedblleft{}To refuse\textquotedblright{}.) Say it once more.
\end{tcolorbox}
